\documentclass[bachelor, och, pract3]{SCWorks}
% параметр - тип обучения - одно из значений:
%    spec     - специальность
%    bachelor - бакалавриат (по умолчанию)
%    master   - магистратура
% параметр - форма обучения - одно из значений:
%    och   - очное (по умолчанию)
%    zaoch - заочное
% параметр - тип работы - одно из значений:
%    referat      - реферат
%    coursework   - курсовая работа (по умолчанию)
%    diploma      - дипломная работа
%    pract2       - отчет по практике, 2 курс, бакалавриат, магистратура
%    pract3       - отчет по практике, 3 курс, бакалавриат
%    pract4       - отчет по практике, 4 курс, бакалавриат
%    practpreddip - отчет по преддипломной практике, 4 курс (= pract3 или nir) 
%    nir          - отчет о научно-исследовательской работе
%    autoref      - автореферат выпускной работы
%    assignment   - задание на выпускную квалификационную работу
%    review       - отзыв руководителя
%    critique     - рецензия на выпускную работу
% параметр - включение шрифта
%    times    - включение шрифта Times New Roman (если установлен)
%               по умолчанию выключен
\usepackage{cmap}
\usepackage[T2A]{fontenc}
\usepackage[utf8]{inputenc}
\usepackage{graphicx}
\usepackage[sort,compress]{cite}
\usepackage{amsmath}
\usepackage{amssymb}
\usepackage{amsthm}
\usepackage{fancyvrb}
\usepackage{longtable}
\usepackage{array}
\usepackage[english,russian]{babel}

% выделение ссылок цветом. чтобы включить- true
\usepackage[colorlinks=false]{hyperref}

\renewcommand{\rmdefault}{cmr} % курсив и полужирный включаются здесь.
\newcommand{\eqdef}{\stackrel {\rm def}{=}}

\newtheorem{lem}{Лемма}

\begin{document}

% Кафедра (в родительном падеже)
\chair{теории функций и стохастического анализа} 
\chairr{кафедра теории функций и стохастического анализа} % для практик!!!!
%\chairr{завод "Тантал"}% для практик!!!!

%Тема работы

\title{Моделирование случайных величин. Вариант 7}

% Курс
\course{2}

% Группа
\group{212}

\department{механико-математического факультета}

% Специальность/направление код - наименование
\napravlenie{01.03.02 "--- Прикладная математика и информатика}
%\napravlenie{38.03.05 "--- Бизнес-информатика}
%\napravlenie{09.04.03 "--- Прикладная информатика}

% Для СТУДЕНТКИ!!!. Для работы студента следующая команда не нужна.
%\studenttitle{Студентки}

% Фамилия, имя, отчество в родительном падеже
\author{Чесакова Максима Евгеньевича}

% Заведующий кафедрой
%\chtitle{д.\,ф.-м.\,н., доцент} % степень, звание
%\chname{С.\,П.\,Сидоров}

%Научный руководитель (для реферата преподаватель проверяющий работу)
\satitle{доцент, к.\,ф.-м.\,н.} %должность, степень, звание
\saname{А.\,Д.\,Луньков} % для 451 гр
%\saname{А.\,К.\,Смирнов} % для 412 гр

% Руководитель практики от организации (только для практики,
% для остальных типов работ не используется)

\patitle{руководитель направления АРМ}
\paname{Д.\,Э.\,Кнутов}

% Семестр (только для практики, для остальных
% типов работ не используется)
\term{5}

% Наименование практики (только для практики, для остальных
% типов работ не используется)
%\practtype{производственная(базовая)}

% Продолжительность практики (количество недель) (только для практики,
% для остальных типов работ не используется)
%\duration{4}

% Даты начала и окончания практики (только для практики, для остальных
% типов работ не используется)
\practStart{29.06.2026}
\practFinish{12.07.2026}

% Год выполнения отчета
\date{2026}

\maketitle

%Включение нумерации рисунков, формул и таблиц по разделам
% (по умолчанию - нумерация сквозная)
% (допускается оба вида нумерации)
%\secNumbering


\tableofcontents

% Раздел "Обозначения и сокращения". Может отсутствовать в работе
%\abbreviations
%\begin{description}
%    \item $|A|$  "--- количество элементов в конечном множестве $A$;
 %   \item $\det B$  "--- определитель матрицы $B$;
 %   \item ИНС "--- Искусственная нейронная сеть;
 %   \item FANN "--- Feedforward Artifitial Neural Network
%\end{description}

% Раздел "Определения". Может отсутствовать в работе
%\definitions

% Раздел "Определения, обозначения и сокращения". Может отсутствовать в работе.
% Если присутствует, то заменяет собой разделы "Обозначения и сокращения" и "Определения"
%\defabbr


% Раздел "Введение"
\intro
%  пример для отчетов по практике для бакалавриата 2, 3 курсов

Целью учебной ознакомительной практики(заменить на соответствующую) <<Правила оформления курсовых и дипломных работ>> (заменить на свою) является создание примера оформления студенческой работы средствами системы \LaTeX.

Поставлена задача оформить документ в соответствии:
\begin{itemize}
    \item со стандартом СТО 06.02-2019 Порядком выполнения, структурой и правилами оформления курсовых работ (проектов)
    и выпускных квалификационных работ, принятых в Саратовском государственном университете в 2019 году;
    \item с правилами оформления титульного листа отчета о прохождении практики в соответствии со стандартом СТО 1.01-2005.
\end{itemize}

При решении поставленных задач (Во время прохождения практики) и при подготовке отчета были использованы источники [1--N]
% N количество источников в списке литературы, для отчета норма не установлена, но от пяти было бы неплохо.
% эту фразу можно не писать, если вы гарантированно расставите ссылки на ВСЕ источники внутри текста отчета.

\section{Стандартный метод моделирования дискретных случайных величин}

\[
	P(\xi = k) = \frac{\sqrt{k+1} - \sqrt{k}}{\sqrt{m}}, \quad 0 \leq k \leq m-1,
\]

\begin{tabular}{cccccccc}
	$\xi$ & 0 & 1 & 2 & 3 & 4 & 5 & 6 \\
	$P$ & 0.377964 & 0.156558 & 0.120131 & 0.101275 & 0.089225 & 0.080666 & 0.074180
\end{tabular}




\section{Моделирование непрерывной случайной величины. Метод обратных функций}


\[
	F_\xi(x) = C(1-\cos(x)), \quad x \in \left[0, \frac{\pi}{4}\right].
\]
\[
C\left(1-\cos\left(\frac{\pi}{4}\right)\right) = 1
\]
\[
C = \frac{1}{1-\cos\left(\frac{\pi}{4}\right)} = 2 + \sqrt{2}
\]
\[
F_\xi(x) = (2 + \sqrt{2})(1 - \cos(x)), \quad x \in \left[0, \frac{\pi}{4}\right]
\]
Обратная функция:
\[
F_\xi^{-1}(y) = \arccos\left(1 - \frac{y}{2 + \sqrt{2}}\right), \quad y \in [0, 1]
\]
Плотность распределения:
\[
f_\xi(x) = (2 + \sqrt{2}) \sin(x), \quad x \in \left[0, \frac{\pi}{4}\right]
\]
Математическое ожидание:
\[
E\xi = \int_{0}^{\frac{\pi}{4}} x f_\xi(x) dx = \int_{0}^{\frac{\pi}{4}} (2 + \sqrt{2}) x \sin(x) dx = 
\]


\subsection{Математическое ожидание}
\[
\mathbb E[\xi]=\int_0^{\pi/4} x\, f_\xi(x)\,dx
= C\int_0^{\pi/4} x\sin x\,dx.
\]
Интегрируем по частям:
\[
\int x\sin x\,dx = -x\cos x+\sin x.
\]
Тогда
\[
\int_0^{\pi/4} x\sin x\,dx
= \left[-x\cos x+\sin x\right]_0^{\pi/4}
= -\frac{\pi}{4}\cdot\frac{\sqrt2}{2}+\frac{\sqrt2}{2}
= \frac{\sqrt2}{2}\left(1-\frac{\pi}{4}\right).
\]
Следовательно,
\[
\mathbb E[\xi]
= C\cdot \frac{\sqrt2}{2}\left(1-\frac{\pi}{4}\right)
= (2+\sqrt2)\cdot\frac{\sqrt2}{2}\left(1-\frac{\pi}{4}\right)
= (\sqrt2+1)\left(1-\frac{\pi}{4}\right).
\]
Итак,
\[
\mathbb E[\xi]=(\sqrt2+1)\left(1-\frac{\pi}{4}\right) \approx 0.51809.
\]

\subsection{Дисперсия}
Используем формулу:
\[
\mathbb D[\xi]=\mathbb E[\xi^2]-(\mathbb E[\xi])^2.
\]
Найдём \(\mathbb E[\xi^2]\):
\[
\mathbb E[\xi^2]=C\int_0^{\pi/4} x^2\sin x\,dx.
\]
Интегрируем по частям дважды (или используем известную первообразную):
\[
\int x^2\sin x\,dx = -x^2\cos x+2x\sin x+2\cos x.
\]
Подставляем пределы:
\[
\int_0^{\pi/4} x^2\sin x\,dx
= \left[-x^2\cos x+2x\sin x+2\cos x\right]_0^{\pi/4}.
\]
Вычисляем:
\[
\begin{aligned}
\left[-x^2\cos x+2x\sin x+2\cos x\right]_0^{\pi/4}
&= -\frac{\pi^2}{16}\cdot\frac{\sqrt2}{2}
+2\cdot\frac{\pi}{4}\cdot\frac{\sqrt2}{2}
+2\cdot\frac{\sqrt2}{2} - (0+0+2) \\
&= \sqrt2\left(1+\frac{\pi}{4}-\frac{\pi^2}{32}\right)-2.
\end{aligned}
\]
Тогда
\[
\mathbb E[\xi^2]=(2+\sqrt2)\left[\sqrt2\left(1+\frac{\pi}{4}-\frac{\pi^2}{32}\right)-2\right].
\]
Раскроем скобки:
\[
\mathbb E[\xi^2]
=2(\sqrt2+1)\left(1+\frac{\pi}{4}-\frac{\pi^2}{32}\right)-4-2\sqrt2.
\]
Теперь квадрат математического ожидания:
\[
(\mathbb E[\xi])^2=(\sqrt2+1)^2\left(1-\frac{\pi}{4}\right)^2
= (3+2\sqrt2)\left(1-\frac{\pi}{4}\right)^2.
\]
Окончательно:

\[
\mathbb D[\xi]
=2(\sqrt2+1)\left(1+\frac{\pi}{4}-\frac{\pi^2}{32}\right)-4-2\sqrt2-(3+2\sqrt2)\left(1-\frac{\pi}{4}\right)^2
.
\]
\[
\mathbb D[\xi] \approx 0.034607.
\]




\section{Моделирование случайных векторов}

\[
f_{\xi, \eta }(x,y) =
\begin{cases}
    \displaystyle\frac{C}{x+y}, \quad x,y \in [1, 2] \times [1, 2] \\
	0, \quad x,y \notin [1, 2] \times [1, 2]
\end{cases}
\]

Найдём константу \(C\), взяв интеграл плотности по области определения:
\[
\int_1^2 \int_1^2 \frac{C}{x+y} \, dy \, dx = 1.
\]
Вычислим интеграл:
\[
\int_1^2 \int_1^2 \frac{C}{x+y} \, dy \, dx =
\]

В результате получаем:
\[
C = \frac{1}{\ln \frac{1024}{729}}.
\]


будем пользоваться методом исключения.

Найдём условную плотность распределения \(\eta\) при \(\xi = x\):
\[
f_{\eta|\xi}(y|x) = \frac{f_{\xi, \eta}(x, y)}{f_\xi(x)}.
\]

Найдём маргинальную плотность \(f_\xi(x)\):
\[
f_\xi(x) = \int_1^2 f_{\xi, \eta}(x, y) \, dy = \int_1^2 \frac{C}{x+y} \, dy = C \ln \frac{x+2}{x+1}.
\]

Найдём маргинальную плотность \(f_\eta(y)\):
\[
f_\eta(y) = \int_1^2 f_{\xi, \eta}(x, y) \, dx = \int_1^2 \frac{C}{x+y} \, dx = C \ln \frac{y+2}{y+1}.
\]

\subsection{Теория из пособия}
\textbf{\S 5. МЕТОД ИСКЛЮЧЕНИЯ (МЕТОД НЕЙМАНА, МЕТОД ОТБОРА)}

Пусть \(g(x) > 0, G = \{(x, y) : 0 \le y \le g(x)\}\) и \(S_G = \int_{-\infty}^{\infty} g(x)dx < \infty\), т. е. площадь области \(G\) конечна. Из курса теории вероятностей известно, что если вектор \((\xi, \eta)\) равномерно распределен в области \(G\), то его плотность распределения вероятностей имеет вид
\[
p(x, y) = \begin{cases} 1/S_G, & (x, y) \in G, \\ 0, & (x, y) \notin G. \end{cases}
\]
Рассмотрим процесс моделирования случайного вектора, равномерно распределенного в области \(G\), т. е. процесс моделирования его координат.

Для плотности распределения вероятностей $p_{\xi}(x)$ первой координаты вектора \((\xi, \eta)\) имеем
\[
p_{\xi}(x) = \int_{-\infty}^{\infty} p(x, y)dy = \int_{0}^{g(x)} \frac{dy}{S_G} = \frac{g(x)}{S_G}.
\]
Так как \(p(x, y) = p_{\xi}(x) \cdot p_{\eta/\xi}(y/x)\), где \(p_{\eta/\xi}(y/x)\) --- условная плотность с. в. \(\eta\) при условии, что \(\xi = x\), то \(p_{\eta/\xi}(y/x) = \frac{p(x, y)}{p_{\xi}(x)} = 1/g(x)\), т. е. при каждом фиксированном \(x\) с. в. \(\eta\) равномерно распределена на отрезке \([0, g(x)]\).

Таким образом, для моделирования случайного вектора, равномерно распределенного в области \(G\), сначала моделируется с. в. \(\xi\) с плотностью \(p_{\xi}(x) = g(x)/S_G\), а после получения ее значения \(\xi = x_0\) моделируется с. в. \(\eta\), равномерно распределенная на отрезке \([0, g(x_0)]\), т. е. \(\eta = \alpha \cdot g(x_0)\), где \(\alpha\) --- значение с. в., равномерно распределенной на отрезке \([0, 1]\); смоделированное значение вектора \((\xi, \eta)\) равно \((x_0, \alpha g(x_0))\).

Изложим теперь метод исключения (см. [6]). Предположим, что нужно моделировать с. в. с плотностью \(p(x)\); пусть \(\xi_1\) --- с. в. c плотностью \(f(x)\), которую мы умеем моделировать, и существует постоянная \(C\) такая, что \\ \(p(x) \le C f(x)\), \(x \in R\). Обозначим через \(G\) и \(H\) области \(\{(x, y) : 0 \le y \le C f(x)\}\) и \(\{(x, y) : 0 \le y \le p(x)\}\) соответственно, а через \((\xi_1, \eta_1)\) и \((\xi, \eta)\) --- случайные векторы, равномерно распределенные в областях \(G\) и \(H\) соответственно. Тогда, как было показано выше, с. в. \(\xi\) имеет плотность распределения \(p(x)\) (в этом случае \(g(x) = p(x)\) и \(S_H = 1\)).

Метод моделирования основан на том, что если \((\xi_1, \eta_1) \in H\), то условное распределение вектора \((\xi_1, \eta_1)\) тоже будет равномерным; действительно, если $D \subset H$, то
\[
P((\xi_1, \eta_1) \in D \mid (\xi_1, \eta_1) \in H) = \frac{P((\xi_1, \eta_1) \in D)}{P((\xi_1, \eta_1) \in H)} = \frac{S_D / S_G}{S_H / S_G} = \frac{S_D}{S_H} = S_D,
\]
(т. е., если случайные точки \((\xi_1, \eta_1)\) равномерно заполняют область \(G\), то те из них, которые попали в \(H\), будут равномерно заполнять \(H\)). Поэтому, если \((\xi_1, \eta_1) \in H\), то полученное для \(\xi_1\) значение принимается в качестве значения с. в. \(\xi\).

Отсюда следует алгоритм моделирования с. в. $\xi$ с плотностью \(p(x)\): моделируем случайный вектор \((\xi_1, \eta_1)\), т. е. моделируем с. в. \(\xi_1\) с плотностью \(f(x)\), и если \(\xi_1 = x_0\), то моделируем с. в. \(\eta_1\), равномерно распределенную на отрезке \([0, C f(x_0)]\); если при этом окажется, что \(\eta_1 < p(x_0)\), то полагаем \(\xi = x_0\), в противном случае моделирование начинается сначала, т. е. моделируются с. в. \(\xi_1\).

Заметим, что плотность \(f(x)\) и величину \(C\) надо выбирать так, чтобы \(C\) было ближе к \(1\); чем \(C\) ближе к \(1\), тем больше доля полезных генераций \((\xi_1, \eta_1)\), ибо вероятность попадания вектора \((\xi_1, \eta_1)\) в \(H\) равна \(1/C\).

Обычно в качестве \(f(x)\) берут ступенчатые или кусочно-линейные функции. Наиболее простой случай получается, когда плотность $p(x)$ определена на конечном промежутке $[a,b]$ и ограничена, т.е. $p(x) \leq M$, $a \leq x \leq b$. Тогда можно положить \(f(x) = 1 / (b-a)\), т. е. с. в. \(\xi_1\) равномерно распределена на отрезке \([a, b]\), и выбрать постоянную $C$ так, что $\displaystyle\frac{C}{(b-a)} \geq M$ (например, $C = \max [1, M(b-a)]$). Поэтому \(\xi_1 = a + \alpha_1(b-a)\) и \(\eta_1 = \alpha_2 M\), где \(\alpha_1, \alpha_2\) --- значения с. в., равномерно распределенной на отрезке \([0, 1]\); в этом случае, если \(\alpha_2 M < p(a + \alpha_1(b-a))\), то полагаем \(\xi = a + \alpha_1(b-a)\).

\textbf{Пример.}

С. в. \(\xi\) имеет плотность \(p(x) = 2\sqrt{1-x^2} / \pi\) для \(-1 \le x \le 1\) и \(p(x) = 0\) при \(x < -1, x > 1\). Так как \(\max p(x) = 2/\pi\), то $\xi_1 = 2 \alpha_1 - 1$ и $\eta_1 = 2 \alpha_2/ \pi$. Тогда положим $\xi = 2 \alpha_1 - 1$, если $2 \alpha_2/ \pi < \frac{2 \sqrt{1 - (2 \alpha - 1)^2}}{\pi}$, т.е. $\alpha_2 < \sqrt{4 \alpha_1 - 4 \alpha_1^2}$.

Таким образом, если $\alpha_2 < 2 \sqrt{\alpha_1(1 - \alpha_1)}$, то $\xi = 2 \alpha_1 - 1$; в противном случае снова разыгрываем пару $(\alpha_1, \alpha_2)$.

Можно указать некоторые обобщения метода отбора (см. [3]). Пусть с. в. \(\xi\) имеет плотность распределения вероятностей \(p(x)\), представимую в виде \(p(x) = C \cdot p_1(x) \cdot g(x)\), где \(C\) --- постоянная, \(p_1(x)\) --- плотность, которую умеем моделировать, и \(0 \le g(x) \le 1\). Тогда значения с. в. \(\xi\) можно получить по следующему алгоритму:
\begin{enumerate}
    \item моделируем с. в. \(\xi\) с плотностью \(p_1(x)\), пусть \(\xi = x_0\);
    \item вычисляем \(g(x_0)\);
    \item моделируем значение \(\alpha\) с. в., равномерно распределенной на отрезке \([0, 1]\);
    \item если \(\alpha < g(x_0)\), то полагаем \(\xi = x_0\); в противном случае моделирование опять начинается с пункта 1.
\end{enumerate}


\noindent \textbf{\S 7. МОДЕЛИРОВАНИЕ НЕПРЕРЫВНЫХ СЛУЧАЙНЫХ ВЕКТОРОВ}

Рассмотрим вначале стандартный метод моделирования векторов. Для моделирования случайного вектора \((\xi, \eta)\) с независимыми координатами можно отдельно моделировать его координаты \(\xi\) и \(\eta\): если \(P_\xi(x)\) и \(P_\eta(y)\) --- плотности распределения вероятностей \(\xi\) и \(\eta\) соответственно, то 
\[ P_\xi(x) = \int_{-\infty}^{\infty} p(x,y)\,dy \quad \text{и} \quad P_\eta(y) = \int_{-\infty}^{\infty} p(u,y)\,du, \]
где \(p(x,y)\) --- плотность распределения случайного вектора \((\xi, \eta)\).

В общем случае, если \(P_\xi(x)\) уже найдена по вышеуказанной формуле, можно вычислить условную плотность \(P_{\eta|\xi}(y|x)\) с.в. $\eta$ при условии $\xi = x$ по формуле \(P_{\eta|\xi}(y|x) = p(x,y) / P_\xi(x)\). Отсюда следующий способ моделирования вектора \((\xi, \eta)\):

1) моделируем с.в. \(\xi\) с плотностью \(P_\xi(x)\), пусть \(\xi = x_0\);

2) находим условную плотность \(P_{\eta|\xi}(y|x_0)\);

3) моделируем с.в. $\eta$ с плотностью \(P_{\eta|\xi}(y|x_0)\), получаем \(\eta = y_0\), тогда \((x_0, y_0)\) есть значение случайного вектора \((\xi, \eta)\).

Указанный способ моделирования очевидно обобщается на вектор любой размерности.

Опишем теперь метод исключения для моделирования случайных векторов. Пусть \(\vec{\xi} = (\xi_1, \xi_2, \dots, \xi_n)\) --- случайный вектор с плотностью распределения вероятностей \(p(x_1, x_2, \dots, x_n)\), причем \(p(x_1, \dots, x_n) \le C\) и каждая с.в.~\(\xi_i\) (\(i=1, \dots, n\)) принимает значения на конечном отрезке \([a_i, b_i]\), \(1 \le i \le n\). Тогда для моделирования случайного вектора \(\vec{\xi}\) можно применить следующий алгоритм:

1) моделируем независимые значения \(\alpha_1, \alpha_2, \dots, \alpha_{n+1}\) с.в., равномерно распределенной на отрезке \([0,1]\);

2) строим вектор \(\vec{z} = (z_1, \dots, z_{n+1})\), где \(z_k, 1 \le k \le n\), есть значение с.в.,  равномерно распределенной на отрезке \([a_k, b_k]\), т.е. \(z_k = a_k + \alpha(b_k - a_k)\), а \(z_{n+1}\) есть значение с.в., равномерно распределенной на \([0, C]\), т.е. $z_{n+1} = \alpha_{n+1}\cdot C$;

3) если \(z_{n+1} < p(z_1, z_2, \dots, z_n)\), то полагаем \(\vec{\xi} = (z_1, \dots, z_n)\), в противном случае значения вектора \(\vec{z} = (z_1, z_2, \dots, z_n)\) отбрасываются и моделирование опять начинается с пункта 1.

Рассмотрим теперь моделирование нормального вектора. Если \(n\)-мер\-ный вектор \(\vec{\xi}\) распределен по многомерному нормальному закону, то \[\vec{\xi}^{~T} = A \cdot \vec{\eta}^{~T} + \vec{m}^{~T},\] где \(A\) --- невырожденная матрица размерности \(n \times n\); \(\vec{\eta}\) --- \(n\)-мерный вектор, все координаты которого независимые одинаково распределенные по нормальному закону \(N(0,1)\) случайные величины; \(\vec{m}\) --- постоянный вектор, при этом \(M\vec{\xi} = \vec{m}\) и корреляционная матрица \(R = AA^T\).

Как известно, многомерное нормальное распределение полностью определяется заданием математического ожидания \(\vec{m}\) и корреляционной матрицы \(R\). Предположим, что матрица \(A\) треугольная, т.е.
\[ A = \begin{pmatrix} 
a_{11} & 0 & 0 & \dots & 0 \\ 
a_{21} & a_{22} & 0 & \dots & 0 \\ 
\dots & \dots & \dots & \dots & \dots \\ 
a_{n1} & a_{n2} & \dots & \dots & a_{nn} 
\end{pmatrix} \]
и покажем, как найти элементы \(a_{ij}, 1 \le i,j \le n\), если известны $R$ и \(\vec{m} = (m_1, m_2, \dots, m_n)\).

Так как \(\xi_1 = a_{11}\eta_1 + m_1\), то \(R_{11} = D\xi_1 = a_{11}^2\), и, следовательно \(a_{11} = \sqrt{R_{11}}\). Аналогично \(\xi_2 = a_{21}\eta_1 + a_{22}\eta_2 + m_2\) и \(R_{12} = M(\xi_1 - m_1)(\xi_2 - m_2) = M[a_{11}\eta_1(a_{21}\eta_1 + a_{22}\eta_2)] = a_{11} a_{21}\), \(R_{22} = D\xi_2 = a_{21}^2 + a_{22}^2\), откуда
\[ a_{21} = \frac{R_{12}}{a_{11}} = \frac{R_{12}}{\sqrt{R_{11}}} \quad \text{и} \quad a_{22} = \sqrt{R_{22} - \frac{R_{12}^2}{R_{11}}}. \]

Общая рекуррентная формула имеет вид
\[ a_{ij} = \frac{R_{ij} - \sum_{k=1}^{j-1} a_{ik} a_{jk}}{\sqrt{R_{jj} - \sum_{k=1}^{j-1}a_{jk}^2}}, \quad 1 \le j \le i \le n, \]
причем считаем, что \(\sum_{k=1}^{0} a_{ik} a_{jk} = 0\). Эта формула легко проверяется рассмотрением величины \(R_{ij}\) для \(j=i\), а затем --- для \(j<i\).


\end{document}
